% ============================================================
% Section 110: 磁化率数值估计（Pauli顺磁 + 迈斯纳效应）
% ============================================================
\section{固体物理：磁化率数值估计}
\label{sec:susceptibility-estimate}

\begin{modelbox}[题目：自由电子气与超导体的磁化率]
给出一个对磁化率 $\chi=M/B$ 的数值估计（精度 $60\%$），其中 $B$ 是实际磁场。

(1) $N=4.7\times10^{22}\,\text{cm}^{-3}$ 的自由电子费米气，在 $T=300\,\text{K}$ 和 $T=350\,\text{K}$ 时；

(2) 在临界温度以下的第一类超导体。
\end{modelbox}

\begin{tipbox}[考点快速分析]
\begin{itemize}
    \item \textbf{Pauli顺磁性}：自由电子气泡利顺磁化率 $\chi=3N\mu_B^2/(2E_F)$，几乎与温度无关
    \item \textbf{CGS vs SI}：CGS 中体积磁化率无量纲；SI 中需乘以 $4\pi$
    \item \textbf{迈斯纳效应}：第一类超导体完全抗磁，$B=0$，$\chi=-1/(4\pi)$（CGS）
\end{itemize}
\end{tipbox}

\begin{derivationbox}[标准解题过程]
\textbf{(1) 自由电子费米气的磁化率（泡利顺磁性）}

自由电子气具有弱顺磁性，来源于费米面附近电子的自旋极化。单位体积泡利顺磁化率（CGS）为
\begin{equation}
\chi_{\text{Pauli}}=\frac{3}{2}\frac{N\mu_B^2}{E_F}.
\end{equation}
其中 $N$ 为电子数密度，$\mu_B=e\hbar/(2m)$ 为玻尔磁子，$E_F$ 为费米能。

\textbf{特点}：泡利顺磁化率与温度几乎无关（$k_BT\ll E_F$ 时），因为只有费米面附近极少数电子参与热激发，经典居里顺磁性被泡利不相容原理强烈抑制。

\textbf{数值计算}（CGS 单位制）：
\begin{align}
N&=4.7\times10^{22}\,\text{cm}^{-3},\\
\hbar&=1.055\times10^{-27}\,\text{erg}\cdot\text{s},\\
m&=9.11\times10^{-28}\,\text{g},\\
\mu_B&=9.27\times10^{-21}\,\text{erg/G}.
\end{align}

费米能（绝对零度近似）：
\begin{equation}
E_F=\frac{\hbar^2}{2m}(3\pi^2N)^{2/3}
\approx\frac{(1.055\times10^{-27})^2}{2\times9.11\times10^{-28}}(3\pi^2\times4.7\times10^{22})^{2/3}
\approx 7.6\times10^{-12}\,\text{erg}\approx4.8\,\text{eV}.
\end{equation}

代入公式：
\begin{equation}
\chi=\frac{3\times4.7\times10^{22}\times(9.27\times10^{-21})^2}{2\times7.6\times10^{-12}}
\approx\frac{1.21\times10^{-17}}{1.52\times10^{-11}}
\approx 8\times10^{-7}.
\end{equation}

由于泡利磁化率不随温度显著变化，在 $T=300\,\text{K}$ 和 $T=350\,\text{K}$ 时，磁化率均约为
\begin{equation}
\boxed{\chi\approx0.8\times10^{-6}\quad(\text{CGS，无量纲}).}
\end{equation}
（SI 单位制中 $\chi_{\text{SI}}=4\pi\chi_{\text{CGS}}\approx1\times10^{-5}$。）

\textbf{(2) 第一类超导体在临界温度以下的磁化率}

第一类超导体在 $T<T_c$、$H<H_c$ 时表现出\textbf{完全抗磁性}（迈斯纳效应）：磁通量完全被排出，内部 $B=0$。

在 CGS 制中 $B=H+4\pi M$，由 $B=0$ 得
\begin{equation}
M=-\frac{H}{4\pi}.
\end{equation}
因此体积磁化率为
\begin{equation}
\boxed{\chi=\frac{M}{H}=-\frac{1}{4\pi}\approx-0.08.}
\end{equation}

对于非理想样品或中间态，有效磁化率绝对值可能小于 $1/(4\pi)$，故
\begin{equation}
\chi\le-\frac{1}{4\pi}.
\end{equation}

\textbf{数量级比较}：超导体的抗磁化率（$\sim-0.1$）比自由电子气的顺磁化率（$\sim10^{-6}$）大五个数量级以上，超导态的抗磁性完全掩盖正常态的弱顺磁性。
\end{derivationbox}

\begin{formulabox}[答案汇总]
\begin{center}
\begin{tabular}{lll}
\toprule
\textbf{系统} & $\boldsymbol{\chi}$ & \textbf{特征} \\
\midrule
自由电子气 & $\approx0.8\times10^{-6}$（与温度无关） & 弱顺磁性（泡利顺磁） \\
第一类超导体 & $\le-1/(4\pi)\approx-0.08$ & 完全抗磁性（迈斯纳效应） \\
\bottomrule
\end{tabular}
\end{center}
\end{formulabox}

\begin{insightbox}[物理图像]
\begin{itemize}
    \item 金属中自由电子的顺磁性远小于经典预期，是泡利不相容原理和费米--狄拉克统计的直接结果。其磁化率与温度无关，仅由费米能级处的态密度决定。
    \item 超导体的完全抗磁性是宏观量子现象，磁化率高达 $-1/(4\pi)$（CGS），是凝聚态物理中最显著的磁现象之一，也是磁悬浮等技术应用的基础。
    \item 注意单位制差异：CGS 中 $\chi$ 无量纲；SI 中超导体的 $\chi=-1$。
\end{itemize}
\end{insightbox}
